\section{Prosedur Algoritma RSA}
\label{sec:prosedur-rsa}

Sistem RSA bekerja melalui tiga tahap utama: pembangkitan kunci, enkripsi, dan dekripsi.
Ketiganya dijalankan dengan rumus yang sangat pendek, tetapi masing-masing menyimpan
syarat yang bila dilanggar akan merusak keamanan atau bahkan membuat pesan tidak dapat
dipulihkan. Subbagian berikut menguraikan ketiga tahap itu beserta syaratnya.

\subsection{Pembangkitan Pasangan Kunci}

Pembangkitan kunci dijalankan satu kali oleh pemilik kunci dan hasilnya dipakai untuk
banyak transaksi. Langkah-langkahnya sebagai berikut.

\begin{enumerate}
    \item Pilih dua bilangan prima besar yang berbeda, $p$ dan $q$.
    \item Hitung modulus $n = p \cdot q$.
    \item Hitung nilai $\varphi(n) = (p - 1)(q - 1)$.
    \item Pilih bilangan bulat $e$ sebagai kunci publik sedemikian sehingga
          $1 < e < \varphi(n)$ dan $\text{PBB}\bigl(e, \varphi(n)\bigr) = 1$.
    \item Hitung kunci privat $d$ sebagai invers modulo dari $e$ terhadap $\varphi(n)$:
          \[ d \equiv e^{-1} \pmod{\varphi(n)} \quad \text{atau} \quad e \cdot d \equiv 1 \pmod{\varphi(n)} . \]
\end{enumerate}

Hasil akhirnya adalah kunci publik $(e, n)$ yang boleh disebarkan dan kunci privat
$(d, n)$ yang harus dijaga. Setelah $d$ diperoleh, nilai $p$, $q$, dan $\varphi(n)$ tidak
diperlukan lagi dan sebaiknya dihapus dari memori --- bila nilai-nilai itu tetap
tersimpan, penyerang yang berhasil membacanya dapat menghitung $d$ seketika.

Nilai $e$ yang paling banyak dipakai adalah $65537 = 2^{16} + 1$, bukan karena angka itu
ajaib melainkan karena tiga alasan praktis: ia cukup besar untuk menghindari serangan
terhadap eksponen kecil, bentuk binernya hanya memuat dua bit bernilai satu sehingga
pemangkatannya cepat, dan ia ganjil serta prima sehingga hampir selalu relatif prima
terhadap $\varphi(n)$. Nilai $e = 3$ juga sah secara matematis, tetapi hanya aman bila
dipakai bersama skema \textit{padding} yang benar seperti dibahas pada
Subbagian~\ref{subsec:padding-oaep}; tanpa \textit{padding}, $e = 3$ dapat dipecahkan
hanya dengan akar pangkat tiga biasa.

\subsection{Syarat Bilangan Prima yang Dipilih}

Tidak setiap pasangan bilangan prima cocok dipakai. Selain harus berbeda dan berukuran
besar, ada empat syarat yang harus dipenuhi agar modulus $n$ benar-benar sukar
difaktorkan:

\begin{itemize}
    \item \textbf{Panjang $p$ dan $q$ hampir sama.} Bila salah satunya jauh lebih pendek,
          penyerang dapat menduga faktor yang kecil dengan mencoba semua bilangan prima
          sampai batas tertentu. Selisih panjang yang dianjurkan tidak lebih dari beberapa
          bit, misalnya dua bilangan prima 1024 bit untuk modulus 2048 bit.
    \item \textbf{Selisih $p - q$ tidak kecil.} Bila $p$ dan $q$ berdekatan, nilai $n$ dapat
          didekati dengan kuadrat salah satunya dan faktorisasinya dapat dicari dengan
          metode Fermat, yang memeriksa $\bigl\lceil\sqrt{n}\bigr\rceil^2 - n$ terhadap
          bilangan kuadrat sempurna. Semakin besar $|p - q|$, semakin banyak langkah yang
          dibutuhkan metode itu.
    \item \textbf{$p - 1$ dan $q - 1$ masing-masing memuat faktor prima besar.} Bila
          keduanya hanya tersusun dari faktor prima kecil, modulusnya rentan terhadap
          algoritma pemfaktoran Pollard $p - 1$ yang bekerja selama batas faktor prima
          terbesarnya masih terjangkau.
    \item \textbf{Bilangan prima dibangkitkan secara acak.} Bilangan prima yang diambil
          dari daftar tetap atau dibangkitkan dengan pembangkit acak yang lemah dapat
          diduga kembali oleh penyerang, betapa pun panjangnya. Pembangkit yang dipakai
          harus memenuhi syarat kriptografis, bukan sekadar pembangkit bilangan
          pseudo-acak biasa.
\end{itemize}

Syarat-syarat tersebut menjelaskan mengapa pembangkitan kunci RSA bukan pekerjaan
sembarang orang: satu pasangan prima yang lemah membuat seluruh sistemnya runtuh,
meskipun panjang kuncinya tampak memadai.

\subsection{Enkripsi dan Pengodean Pesan}

Rumus enkripsinya sendiri hanya satu baris. Yang menuntut kehati-hatian adalah mengubah
pesan berbentuk teks menjadi bilangan bulat $m$ yang memenuhi syarat panjang.

\begin{enumerate}
    \item \textit{Preprocessing}: pesan dikodekan menjadi bilangan bulat $m$. Bila
          hasil pengodeannya lebih besar atau sama dengan $n$, pesan harus dipecah menjadi
          blok-blok $m_1, m_2, \dots, m_t$ dengan syarat
          \begin{equation}
          0 \le m_i \le n - 1 \quad \text{untuk setiap } i = 1, 2, \dots, t .
          \label{eq:syarat-blok}
          \end{equation}
    \item Hitung cipherteks setiap blok dengan kunci publik $(e, n)$:
          \begin{equation}
          c_i = m_i^{\,e} \bmod n .
          \label{eq:enkripsi-blok}
          \end{equation}
\end{enumerate}

Penentuan panjang blok mengikuti batas pada Persamaan~\eqref{eq:syarat-blok}. Bila pesan
dikodekan menjadi $k$ digit desimal per blok, maka blok terbesar yang mungkin adalah
$10^k - 1$, sehingga harus berlaku $10^k - 1 \le n - 1$, yaitu

\begin{equation}
10^k \le n .
\label{eq:panjang-blok}
\end{equation}

Untuk modulus $n = 3337$ yang terdiri atas empat digit, syarat $10^k \le n$ baru
terpenuhi pada $k = 3$, sebab $10^4 = 10000$ melebihi $3337$. Meskipun demikian, pada
contoh terhitung nanti tetap dipakai blok empat digit. Alasannya, syarat yang menentukan
bukan panjang $n$ dalam digit melainkan nilai terbesar yang benar-benar muncul pada pesan.
Kode dua huruf paling besar pada pengodean yang dipakai adalah $2525$, dan
$2525 < 3337$ sehingga Persamaan~\eqref{eq:syarat-blok} tetap terpenuhi walaupun
kapasitas teoretis blok empat digit belum seluruhnya terpakai. Pilihan panjang blok
karena itu harus diperiksa terhadap pesan yang sesungguhnya dikirim, bukan sekadar
terhadap panjang modulus.

Satu hal kecil yang mudah terlewat adalah \textbf{penjagaan angka nol di depan}. Pada
contoh berkode huruf, huruf A dikodekan sebagai $00$. Bila blok $0704$ dituliskan tanpa
nol di depan menjadi $704$, hasil enkripsinya berbeda dan, yang lebih berbahaya, jumlah
digitnya berkurang sehingga batas blok berikutnya bergeser. Penerima akan memulihkan
deretan angka yang tidak sejajar lagi dengan pengodean asalnya. Karena itu panjang setiap
blok harus tetap sama --- dalam contoh ini selalu empat digit --- dan penulisan kembali
nol di depan termasuk bagian dari prosedur, bukan sekadar kerapian penulisan.

\subsection{Dekripsi dan Pengodean Kembali}

\begin{enumerate}
    \item Hitung kembali setiap blok plainteks dengan kunci privat $(d, n)$:
          \begin{equation}
          m_i = c_i^{\,d} \bmod n .
          \label{eq:dekripsi-blok}
          \end{equation}
    \item \textit{Postprocessing}: rangkaikan kembali seluruh $m_i$ menjadi deretan angka
          yang sama panjangnya seperti pada pengodean semula, lalu terjemahkan kembali
          menjadi huruf.
\end{enumerate}

Kebenaran Persamaan~\eqref{eq:dekripsi-blok} telah ditunjukkan pada
Subbagian~\ref{sec:landasan-rsa}. Yang perlu diperhatikan di sini adalah bahwa dekripsi
memerlukan $n$ yang sama dengan yang dipakai pada enkripsi. Modulus $n$ muncul pada kedua
kunci justru karena alasan itu: ia menyatakan ruang tempat seluruh perhitungan modulo
berlangsung, sehingga tanpa modulus yang sama penerima tidak akan memperoleh hasil yang
konsisten.

Langkah-langkah pada ketiga tahap itu dapat dibaca sebagai satu alur pada
Gambar~\ref{fig:alur-rsa}. Perhatikan bahwa kunci publik mengalir menuju kotak enkripsi
sebagai tanda bahwa siapa pun dapat mengenkripsi, sedangkan kunci privat hanya menuju
kotak dekripsi sebagai tanda bahwa hanya pemiliknya yang dapat membuka pesan.

% Alur algoritma RSA: pembangkitan kunci, enkripsi dengan kunci publik,
% dekripsi dengan kunci privat.
% Dirujuk dari section-03.tex dengan % Alur algoritma RSA: pembangkitan kunci, enkripsi dengan kunci publik,
% dekripsi dengan kunci privat.
% Dirujuk dari section-03.tex dengan % Alur algoritma RSA: pembangkitan kunci, enkripsi dengan kunci publik,
% dekripsi dengan kunci privat.
% Dirujuk dari section-03.tex dengan \input{figures/alur-rsa}.
\begin{figure}[htbp]
    \centering
    \begin{tikzpicture}[
        kotak/.style={draw, rounded corners, align=center, font=\small, inner sep=5pt},
        panah/.style={-{Latex[length=2.4mm]}, thick},
        kunci/.style={-{Latex[length=2.4mm]}, thick, dashed, gray!75},
    ]
        \node[kotak, fill=gray!10, minimum width=8.4cm] (gen) at (2.6,0)
            {Pembangkitan kunci: pilih $p$ dan $q$; hitung $n = pq$ dan $\phi(n) = (p-1)(q-1)$;\\
             pilih $e$ lalu hitung $d \equiv e^{-1} \pmod{\phi(n)}$};
        \node[kotak, fill=blue!6, minimum width=3.0cm] (pub)  at (0.4,-1.8) {Kunci publik $(e, n)$};
        \node[kotak, fill=red!6,  minimum width=3.0cm] (priv) at (5.6,-1.8) {Kunci privat $(d, n)$};

        \draw[panah] (gen.south) -- ++(0,-0.4) -| (pub.north);
        \draw[panah] (gen.south) -- ++(0,-0.4) -| (priv.north);

        \node[kotak, fill=green!10,  minimum width=2.4cm] (m)   at (-2.6,-3.7) {Pesan $m$};
        \node[kotak, fill=orange!12, minimum width=3.0cm] (enc) at (0.4,-3.7)  {$c = m^{e} \bmod n$};
        \node[kotak, fill=green!10,  minimum width=1.4cm] (c)   at (2.9,-3.7)  {$c$};
        \node[kotak, fill=orange!12, minimum width=3.0cm] (dec) at (5.6,-3.7)  {$m = c^{d} \bmod n$};
        \node[kotak, fill=green!10,  minimum width=2.4cm] (m2)  at (8.3,-3.7)  {Pesan $m$};

        \draw[kunci] (pub.south)  -- (enc.north);
        \draw[kunci] (priv.south) -- (dec.north);

        \draw[panah] (m) -- (enc);
        \draw[panah] (enc) -- (c);
        \draw[panah] (c) -- (dec);
        \draw[panah] (dec) -- (m2);
    \end{tikzpicture}
    \caption{Alur algoritma RSA: enkripsi memakai kunci publik penerima, dekripsi hanya dapat dilakukan dengan kunci privatnya}
    \label{fig:alur-rsa}
\end{figure}
.
\begin{figure}[htbp]
    \centering
    \begin{tikzpicture}[
        kotak/.style={draw, rounded corners, align=center, font=\small, inner sep=5pt},
        panah/.style={-{Latex[length=2.4mm]}, thick},
        kunci/.style={-{Latex[length=2.4mm]}, thick, dashed, gray!75},
    ]
        \node[kotak, fill=gray!10, minimum width=8.4cm] (gen) at (2.6,0)
            {Pembangkitan kunci: pilih $p$ dan $q$; hitung $n = pq$ dan $\phi(n) = (p-1)(q-1)$;\\
             pilih $e$ lalu hitung $d \equiv e^{-1} \pmod{\phi(n)}$};
        \node[kotak, fill=blue!6, minimum width=3.0cm] (pub)  at (0.4,-1.8) {Kunci publik $(e, n)$};
        \node[kotak, fill=red!6,  minimum width=3.0cm] (priv) at (5.6,-1.8) {Kunci privat $(d, n)$};

        \draw[panah] (gen.south) -- ++(0,-0.4) -| (pub.north);
        \draw[panah] (gen.south) -- ++(0,-0.4) -| (priv.north);

        \node[kotak, fill=green!10,  minimum width=2.4cm] (m)   at (-2.6,-3.7) {Pesan $m$};
        \node[kotak, fill=orange!12, minimum width=3.0cm] (enc) at (0.4,-3.7)  {$c = m^{e} \bmod n$};
        \node[kotak, fill=green!10,  minimum width=1.4cm] (c)   at (2.9,-3.7)  {$c$};
        \node[kotak, fill=orange!12, minimum width=3.0cm] (dec) at (5.6,-3.7)  {$m = c^{d} \bmod n$};
        \node[kotak, fill=green!10,  minimum width=2.4cm] (m2)  at (8.3,-3.7)  {Pesan $m$};

        \draw[kunci] (pub.south)  -- (enc.north);
        \draw[kunci] (priv.south) -- (dec.north);

        \draw[panah] (m) -- (enc);
        \draw[panah] (enc) -- (c);
        \draw[panah] (c) -- (dec);
        \draw[panah] (dec) -- (m2);
    \end{tikzpicture}
    \caption{Alur algoritma RSA: enkripsi memakai kunci publik penerima, dekripsi hanya dapat dilakukan dengan kunci privatnya}
    \label{fig:alur-rsa}
\end{figure}
.
\begin{figure}[htbp]
    \centering
    \begin{tikzpicture}[
        kotak/.style={draw, rounded corners, align=center, font=\small, inner sep=5pt},
        panah/.style={-{Latex[length=2.4mm]}, thick},
        kunci/.style={-{Latex[length=2.4mm]}, thick, dashed, gray!75},
    ]
        \node[kotak, fill=gray!10, minimum width=8.4cm] (gen) at (2.6,0)
            {Pembangkitan kunci: pilih $p$ dan $q$; hitung $n = pq$ dan $\phi(n) = (p-1)(q-1)$;\\
             pilih $e$ lalu hitung $d \equiv e^{-1} \pmod{\phi(n)}$};
        \node[kotak, fill=blue!6, minimum width=3.0cm] (pub)  at (0.4,-1.8) {Kunci publik $(e, n)$};
        \node[kotak, fill=red!6,  minimum width=3.0cm] (priv) at (5.6,-1.8) {Kunci privat $(d, n)$};

        \draw[panah] (gen.south) -- ++(0,-0.4) -| (pub.north);
        \draw[panah] (gen.south) -- ++(0,-0.4) -| (priv.north);

        \node[kotak, fill=green!10,  minimum width=2.4cm] (m)   at (-2.6,-3.7) {Pesan $m$};
        \node[kotak, fill=orange!12, minimum width=3.0cm] (enc) at (0.4,-3.7)  {$c = m^{e} \bmod n$};
        \node[kotak, fill=green!10,  minimum width=1.4cm] (c)   at (2.9,-3.7)  {$c$};
        \node[kotak, fill=orange!12, minimum width=3.0cm] (dec) at (5.6,-3.7)  {$m = c^{d} \bmod n$};
        \node[kotak, fill=green!10,  minimum width=2.4cm] (m2)  at (8.3,-3.7)  {Pesan $m$};

        \draw[kunci] (pub.south)  -- (enc.north);
        \draw[kunci] (priv.south) -- (dec.north);

        \draw[panah] (m) -- (enc);
        \draw[panah] (enc) -- (c);
        \draw[panah] (c) -- (dec);
        \draw[panah] (dec) -- (m2);
    \end{tikzpicture}
    \caption{Alur algoritma RSA: enkripsi memakai kunci publik penerima, dekripsi hanya dapat dilakukan dengan kunci privatnya}
    \label{fig:alur-rsa}
\end{figure}

